\chapter{Pricing of Derivatives}

\section{The General Pricing Formula}

In a financial market, a derivative (or contingent claim) is a contract whose payoff depends on the behavior of the underlying primary assets. Mathematically, a European derivative expiring at maturity $T$ is defined as an $\mathcal{F}_T$-measurable random variable $H$. The $\mathcal{F}_T$-measurability ensures that the payoff $H$ is completely determined by the information available in the market up to time $T$.

In the context of our binomial model with one risky asset, the payoff of any derivative can be expressed as a deterministic function $h$ of the entire price path:
\begin{equation}
    H = h(S_0, S_1, \dots, S_T).
\end{equation}

According to the fundamental precepts of risk-neutral valuation, in an arbitrage-free and complete market, the unique no-arbitrage price of the derivative at any time $t \in \{0, \dots, T\}$, denoted by $V_t$, is given by the conditional expectation of the discounted payoff under the unique equivalent martingale measure $\mathbb{Q}$:
\begin{equation}
    V_t = \mathbb{E}^\mathbb{Q} \left[ \frac{H}{(1+r)^{T-t}} \;\middle|\; \mathcal{F}_t \right].
\end{equation}

Since the filtration $\mathcal{F}_t$ is generated by the stock prices up to time $t$, we intuitively expect the price $V_t$ to depend only on the observed path $(S_0, \dots, S_t)$. The following theorem formalizes this using the properties of the binomial model.

\begin{theorem}[Pricing Function for Path-Dependent Derivatives]
    Under the binomial model, the value of the replicating portfolio for a derivative with payoff $H = h(S_0, S_1, \dots, S_T)$ at time $t$ is given by:
    \begin{equation}
        V_t = v_t(S_0, S_1, \dots, S_t), \quad t = 0, \dots, T,
    \end{equation}
    where the deterministic pricing function $v_t : \mathbb{R}^{t+1} \to \mathbb{R}$ is defined explicitly by replacing the unknown future returns with their random variables under $\mathbb{Q}$:
    \begin{align}
        v_t(x_0, \dots, x_t) &:= \frac{1}{(1+r)^{T-t}} \times \nonumber \\
        &\quad \mathbb{E}^\mathbb{Q} \left[ h\left(x_0, \dots, x_t, x_t(1+R_{t+1}), \dots, x_t \prod_{k=t+1}^T (1+R_k)\right) \right]. \label{eq:general_vt}
    \end{align}
\end{theorem}

\begin{proof}
    By the risk-neutral pricing formula, $V_t = \frac{1}{(1+r)^{T-t}} \mathbb{E}^\mathbb{Q}[h(S_0, \dots, S_T) \mid \mathcal{F}_t]$. 
    We can rewrite the future prices $S_{t+1}, \dots, S_T$ in terms of the time-$t$ price $S_t$ and the future one-period returns $R_{t+1}, \dots, R_T$:
    \begin{equation}
        S_k = S_t \prod_{i=t+1}^k (1+R_i), \quad \text{for } k > t.
    \end{equation}
    Therefore, the payoff $H$ becomes:
    \begin{equation}
        H = h\left(S_0, \dots, S_t, S_t(1+R_{t+1}), \dots, S_t \prod_{k=t+1}^T (1+R_k)\right).
    \end{equation}
    Observe that $(S_0, \dots, S_t)$ are structurally $\mathcal{F}_t$-measurable (they are known constants given the information at time $t$), while the future returns $(R_{t+1}, \dots, R_T)$ are independent of $\mathcal{F}_t$ under the measure $\mathbb{Q}$ (since returns are i.i.d. under $\mathbb{Q}$). 
    
    By the Independence Lemma for conditional expectations, if we have a function $g(X, Y)$ where $X$ is $\mathcal{F}_t$-measurable and $Y$ is independent of $\mathcal{F}_t$, then $\mathbb{E}[g(X,Y) \mid \mathcal{F}_t] = \phi(X)$ where $\phi(x) = \mathbb{E}[g(x, Y)]$.
    Applying this powerful lemma, we substitute dummy variables $(x_0, \dots, x_t)$ for the known path $(S_0, \dots, S_t)$ and take the unconditional expectation over the remaining independent random variables to recover the exact deterministic function $v_t$ provided in equation \eqref{eq:general_vt}.
\end{proof}

\section{Pricing Vanilla Options}

A European derivative is called path-independent or a "vanilla" option if its payoff depends solely on the terminal asset price $S_T$, rather than the entire historical trajectory. In this case, $H = h(S_T)$ for some function $h : \mathbb{R} \to \mathbb{R}$ (for instance, $h(x) = \max(x-K, 0)$ for a European Call).

By the Markov property of the binomial price process under $\mathbb{Q}$, the price of a vanilla option at time $t$ depends strictly upon the current price $S_t$, allowing us to dramatically simplify $v_t(S_0, \dots, S_t)$ to merely $v_t(S_t)$.

\begin{proposition}[Cox-Ross-Rubinstein Formula for Time 0 Pricing]
    Let $H = h(S_T)$ be a path-independent payoff. Its fundamental no-arbitrage price at inception $t=0$ is structurally quantified by:
    \begin{equation}
        v_0(S_0) = \frac{1}{(1+r)^T} \sum_{k=0}^T \binom{T}{k} q^k (1-q)^{T-k} h\left(S_0(1+u)^k(1+d)^{T-k}\right).
    \end{equation}
\end{proposition}

\begin{proof}
    At time $t=0$, the risk-neutral formula declares $v_0(S_0) = \frac{1}{(1+r)^T} \mathbb{E}^\mathbb{Q}[h(S_T)]$. Because $S_T = S_0 \prod_{i=1}^T (1+R_i)$, the terminal price depends entirely upon the total number of "up" movements that unfold dynamically over the $T$ periods. 
    
    Let $X$ denote the absolute number of "up" movements (where $R_i = u$). Since the returns $R_i$ are independent under $\mathbb{Q}$ and identically yield $u$ with strict probability $q$, the variable $X$ intrinsically follows a Binomial distribution: $X \sim \mathcal{B}(T, q)$.
    If $X = k$, then there are exactly $T-k$ "down" movements, driving the terminal price flawlessly to $S_0(1+u)^k(1+d)^{T-k}$.
    Taking the rigorous mathematical expectation over the explicit probability mass function of the binomial distribution directly yields the weighted summation formula above.
\end{proof}

This explicit formula is highly practically efficient, requiring only $\mathcal{O}(T)$ operations. However, to track the price path dynamically through the tree or to price American options, backward induction stands as the preferred algorithmic strategy.

\begin{proposition}[Backward Induction Algorithm]
    For a generalized derivative, the deterministic pricing function $v_t$ thoroughly satisfies the recursive one-step relation functionally bridging time $t$ to $t+1$:
    \begin{align} \label{eq:backward_induction}
        v_t(S_0, \dots, S_t) &= \frac{1}{1+r} \Big[ q v_{t+1}(S_0, \dots, S_t, S_t(1+u)) \nonumber \\
        &\quad + (1-q) v_{t+1}(S_0, \dots, S_t, S_t(1+d)) \Big]
    \end{align}
    subject to the strict terminal boundary condition at maturity $T$:
    \begin{equation}
        v_T(S_0, \dots, S_T) = h(S_0, \dots, S_T).
    \end{equation}
\end{proposition}

\begin{proof}
    This powerful relation flows seamlessly from the innate tower property of conditional expectations. Specifically:
    \begin{equation}
        V_t = \mathbb{E}^\mathbb{Q} \left[ \frac{V_{t+1}}{1+r} \;\middle|\; \mathcal{F}_t \right].
    \end{equation}
    Since $V_{t+1} = v_{t+1}(S_0, \dots, S_t, S_{t+1})$ and $S_{t+1} = S_t(1+R_{t+1})$, we condition entirely upon the $\mathcal{F}_t$-measurable sequence $(S_0, \dots, S_t)$. The single element rigidly remaining uncertain conditionally represents $R_{t+1}$. Under the structural paradigm dictated by $\mathbb{Q}$, $R_{t+1}$ equals $u$ with explicit probability $q$ mapping the future path toward state $(S_0, \dots, S_t, S_t(1+u))$, or rigidly falls to $d$ with probability $(1-q)$ navigating heavily to $(S_0, \dots, S_t, S_t(1+d))$. Taking the expectation over this strictly isolated one-period jump identically produces the recursive algorithm in \eqref{eq:backward_induction}.
\end{proof}
